The apparent magnitude of star (a) is $m_a$, of star (b) is $m_b$ and that of
the system as a whole is $m_{a+b}$. The corresponding apparent brightnesses are
$\ell_a$, $\ell_b$ and $\ell_{a+b} = \ell_a + \ell_b$.

For star (a):
\[
m_{a+b} - m_a = -2.5 \log\!\left(\frac{\ell_a + \ell_b}{\ell_a}\right)
\]
and because $\dfrac{\ell_b}{\ell_a} = \dfrac{1}{2}$, we get
$m_a = m_{a+b} + 2.5 \log\left(1 + \tfrac{1}{2}\right)$ or
$m_a = 5 + 2.5 \log (3/2)$ and finally $m_a = 5.44$~mag. \hfill(5 point)

Similarly for star (b):
\[
m_{a+b} - m_b = -2.5 \log\!\left(\frac{\ell_a + \ell_b}{\ell_b}\right)
\]
and because $\dfrac{\ell_a}{\ell_b} = 2$, we get
$m_b = m_{a+b} + 2.5 \log (3)$ or
$m_b = 5 + 2.5 \log (3)$ and finally $m_b = 6.19$~mag. \hfill(5 Point)
